\documentclass[aspectratio=169,10pt,english]{beamer}

\input{../../../_preambulo_beamer.tex}
\input{../../../_comandos_beamer.tex}

\selectlanguage{english}

\title{MA1001B - Statistical Modeling for Decision Making}
\subtitle{Lesson 15: Binomial Applications and Business Risk}
\author{}
\institute{Tecnol\'ogico de Monterrey}
\date{}

\begin{document}

\begin{frame}[plain]
  \vspace{1.2cm}
  {\Large\bfseries MA1001B - Statistical Modeling for Decision Making\par}
  \vspace{0.7cm}
  {\large Lesson 15: Binomial Applications and Business Risk\par}
  \vspace{0.7cm}
  {\color{TecRojo}\rule{\textwidth}{1.2pt}}
  \vspace{0.8cm}

  MA1001B Faculty\\[0.3cm]
  Term\\[0.3cm]
  Tecnol\'ogico de Monterrey
\end{frame}

\begin{frame}{Two Policies, One Mean}
  \begin{alertblock}{Guiding question}
    If two binomial inspection policies both have $np=2$, do they have the
    same chance of many defectives?
  \end{alertblock}

  \vspace{0.6em}
  By the end of this lesson, you can:
  \begin{itemize}
    \item state Policy A ($n=50$, $p=0.04$) and Policy B ($n=20$, $p=0.10$);
    \item compute $P(X\ge 3)$ for each;
    \item compare $P(X=0)$ and $P(X\ge 8)$;
    \item explain why matching $np$ does not match tail risk.
  \end{itemize}

  \vfill
  \scriptsize All policies used today are synthetic. No real company data are used.
\end{frame}

\begin{frame}{Same $np$, Different $n$ and $p$}
  \small
  \begin{center}
    \begin{tabular}{lccc}
      \toprule
      \textbf{Policy} & $n$ & $p$ & $np$ \\
      \midrule
      A & 50 & 0.04 & 2 \\
      B & 20 & 0.10 & 2 \\
      \bottomrule
    \end{tabular}
  \end{center}

  \vspace{0.5em}
  Variances: $\Var(X_A)=1.920000$ and $\Var(X_B)=1.800000$. Policy A has more
  room for extreme counts because $n$ is larger.
\end{frame}

\begin{frame}[fragile]{Step 1: Two Policies}
  \begin{columns}[T]
    \begin{column}{0.42\textwidth}
      \small
      \begin{lstlisting}
n_a, p_a = 50, 0.04
n_b, p_b = 20, 0.10
mean_a = n_a * p_a
mean_b = n_b * p_b
      \end{lstlisting}

      \begin{block}{Verified output}
        $np_A=2.000000$\\
        $np_B=2.000000$\\
        $\Var_A=1.920000$\\
        $\Var_B=1.800000$
      \end{block}
    \end{column}
    \begin{column}{0.58\textwidth}
      \centering
      \includegraphics[width=\textwidth]{../figures/binomial_risk_01_two_policies.png}
      \vspace{0.2em}
      \scriptsize Policy ingredients from Step 1
    \end{column}
  \end{columns}
\end{frame}

\begin{frame}{$P(X\ge 3)$ Is Not the Whole Risk Story}
  \small
  \[
    P(X_A\ge 3)=0.323286,\qquad P(X_B\ge 3)=0.323073.
  \]
  Those near-tail values almost match. The far tail does not:
  \[
    P(X_A\ge 8)=0.000781,\qquad P(X_B\ge 8)=0.000416.
  \]
  The ratio is $1.880028$. Policy A also has the larger clean-sample chance
  $P(X=0)=0.129886$ versus $0.121577$.
\end{frame}

\begin{frame}[fragile]{Step 2: Tail Probabilities}
  \begin{columns}[T]
    \begin{column}{0.40\textwidth}
      \small
      \begin{lstlisting}
p_ge3 = stats.binom.sf(2, n, p)
p_ge8 = stats.binom.sf(7, n, p)
      \end{lstlisting}

      \begin{block}{Verified output}
        $P(X\ge 3)\approx 0.323$ both\\
        $P_A(X\ge 8)=0.000781$\\
        $P_B(X\ge 8)=0.000416$
      \end{block}
    \end{column}
    \begin{column}{0.60\textwidth}
      \centering
      \includegraphics[width=\textwidth]{../figures/binomial_risk_02_tail_probabilities.png}
      \vspace{0.2em}
      \scriptsize $P(X=0)$ and $P(X\ge 8)$ from Step 2
    \end{column}
  \end{columns}
\end{frame}

\begin{frame}{Step 3: Matching $np$ Hides a Heavier Far Tail}
  \begin{columns}[T]
    \begin{column}{0.42\textwidth}
      \small
      Overlay the two PMFs. The $x\ge 8$ region is almost twice as likely under
      Policy A.

      \vspace{0.4em}
      \textbf{Limit:} same mean $np$ can hide different tail probabilities.
      Lesson 16 draws without replacement, so trials are no longer independent.
    \end{column}
    \begin{column}{0.58\textwidth}
      \centering
      \includegraphics[width=\textwidth]{../figures/binomial_risk_03_same_mean_different_tails.png}
      \vspace{0.2em}
      \scriptsize PMF overlay from Step 3
    \end{column}
  \end{columns}
\end{frame}

\begin{frame}{Concept Check}
  \begin{enumerate}
    \item Why do both policies have mean $2$?
    \item Compute $P(X\ge 3)$ for Policy A from the verified output.
    \item Why can $P(X\ge 3)$ match when $P(X\ge 8)$ does not?
    \item Which policy has the heavier far tail?
  \end{enumerate}

  \vspace{0.6em}
  \begin{block}{Connection to Lesson 16}
    The session can continue directly into the hypergeometric model for
    sampling without replacement.
  \end{block}

  \vfill
  \scriptsize Source: OpenStax, \textit{Introductory Business Statistics 2e},
  Section 4.2. CC BY-NC-SA.
\end{frame}

\appendix



\end{document}
